\chapter{Products, partial sums and measurement boundaries}
\MARSChapterHeading{Products and measurement boundaries}
\label{revision-lab:measurement}

Import the development lesson \texttt{022-measuring-programs} and choose Continue
in Learn. Read the main chapter's arithmetic contract and measurement scopes first.
Its complete \texttt{matrix-measure.asm} is also supplied in
\texttt{revision/labs/files}; both public matrix entries are worked functions.
The matrix-vector interface is your independent exercise.

Learn follows the ordinary dot product, a wide sum, the blocked partial sums and
two larger cache observations through nineteen Debug stops. Nine predictions and
ordinary Run of five complete sources lead to the independent checks. The
alternate files are in \texttt{worked}: \texttt{matrix-wide.asm} selects the
positive $2^{31}$ example, \texttt{matrix-tiled.asm} selects method one, and
\texttt{cache-ordinary.asm}/\texttt{cache-tiled.asm} contain deterministic
sixteen-by-sixteen inputs with tile edge four. Keep these separate from the default
source. Without Learn, use the same boundaries below and make disposable copies
for changes.

\section{Predict a product and its stored width}

Derive all four default output values from rows and columns, then write each as
two little-endian words. Start a fresh Debug session for \texttt{matrix-measure.asm}
and use one line breakpoint at a time. At the first \texttt{naive\_dot}, the
reduction index and accumulator are zero. The input cursors name the first row
and column pair. Resume to the next visit: the first product, seven, is in
\texttt{\$t3}, while the left cursor advanced four bytes and the right cursor
eight. Those strides follow different input layouts, even though both elements
occupy one word.

Move to \texttt{naive\_store}. All three pairs have contributed to
\texttt{\$t3}=58; the output store has not executed yet. An output cell becomes
complete only when that full sum is stored. Learn's source highlight identifies
the next requested stop; Debug's selected stack frame identifies the actual pause.
Continue to \texttt{matmul\_done}. The function has returned,
but the caller has not stored its count in \texttt{completed}. Compare the four
doublewords and \texttt{\$v0}=4 with your calculation. Continue to
\texttt{results\_ready} to inspect the published count before exit.

Finish or stop that execution before opening \texttt{worked/matrix-wide.asm}.
Its shape is $1,2,1,2$ and both active inputs are $[-32768,-32768]$. Predict the
low and high output words for positive $2^{31}$. Break at \texttt{naive\_store}
to inspect the complete positive accumulator before its store, then compare the
doubleword at \texttt{matmul\_done}. At \texttt{results\_ready}, the caller
has published count one before selecting exit. In a disposable copy, replace doubleword accumulation with word accumulation. Find the
first lost upper bit rather than diagnosing the result as a problem with negative
input values. Restore the original before taking reference measurements.

\begin{figure}[htbp]
\centering
\begin{tikzpicture}[alt={An actual native Memory view at matmul_done, cropped to its address control and first two rows. At address hexadecimal 10010900, the first eight bytes are 00 00 00 80 00 00 00 00. Native boxes mark the low word 80000000 and high word zero. Together they represent positive 2147483648. The next eight bytes retain the output sentinel a5 a5 5a 5a a5 a5 5a 5a.}]
\node[anchor=south west,inner sep=0] (shot) at (0,0) {\MARSRevisionCrop{revision/captures/measurement/matrix-wide.png}{0.97\linewidth}{0.032813}{0.311112}{0.142187}{0.524444}};
\begin{scope}[shift={(shot.south west)},x={($(shot.south east)-(shot.south west)$)},y={($(shot.north west)-(shot.south west)$)}]
\draw[black,line width=.7pt] (.60,.16) rectangle (.69,.30);
\draw[black,line width=.7pt] (.69,.16) rectangle (.78,.30);
\end{scope}
\end{tikzpicture}
\caption{The sum of two products $(-32768)(-32768)$ is positive $2^{31}$. In the native Memory view, the first boxed word is the low half and the second is the high half. Their combined doubleword is valid; interpreting only the low word as signed would give the wrong sign. The following doubleword remains a sentinel.}
\label{fig:revision-measurement-width}
\end{figure}


Repeat with $K=0$ and a nonzero sentinel in output. Every active doubleword becomes
zero without input reads. With $M=0$ or $N=0$, no matrix bytes are touched. Explain
why these two empty cases have different output obligations even though neither
performs a multiplication.

\section{Inspect the handoff between reduction blocks}

Start a fresh Debug session for \texttt{worked/matrix-tiled.asm}; only method
selection differs from the default source. At the first \texttt{block\_reduction}
visit, \texttt{\$t0} is zero and the product has been initialized. At its second
visit, \texttt{\$t0} is two; inspect partial outputs $[25,28,73,82]$ in row-major
order. Use a conditional breakpoint on \texttt{\$t0 == 2} or continue past the
first visit. The final reduction contributes the remaining $q=2$ terms.

Move to \texttt{reduction\_bound} to inspect the clipped reduction end in
\texttt{\$a2}. It becomes three, not four,
when $L=2$, $K=3$ and $T=2$. Source pointers begin at that reduction offset;
restarting them at column zero would add earlier terms twice. Replacing the loaded
partial sum with zero would keep only the final block. Try each fault separately
and identify the first incorrect output transition.

Now use $M=5,K=7,N=3,T=3$. Derive the final row, column and reduction bounds on
paper. The final output tile is two by three, while its last reduction interval
has length one. A full-tile loop bound on any of these axes can access a mapped
but unselected caller byte. The checks inspect footprints as well as final values.

\section{Reproduce cold, warm and reset observations}

Run the measurement tool from the repository root, using JDK 25 and a new output
directory:
\begin{verbatim}
python3 tools/measure_revision_matmul.py \
  --java /path/to/jdk-25/bin/java \
  --cli /path/to/mars-cli.jar \
  --output /new/measurement-directory
\end{verbatim}
The tool checks products, sentinel suffixes and exact reference traffic accounting.
It records raw output, source/CLI/tool hashes and a structured measurement receipt.
Its reference counts belong to the supplied source. For a different implementation,
retain its correctness checks and establish a separate observation; a new loop
need not reproduce these exact counts.

Choose the sixteen-by-sixteen, tile-four rows. Compare cold function, second call,
reset replay and launch scopes. Explain why the warm row is a counter difference,
why reset reproduces the cold values, and why launch has one more data notice.
Do not infer that a second IDE launch inherits a prior session's modeled cache.
The warm case is a second call inside the same execution session.

Account for 36 stack notices and four descriptor notices before counting matrix
traffic. For each method, derive the input and product notice totals. A single
doubleword store contributes two native notices here. Compare those totals with
completed instructions and explain why their units differ.

\section{Connect the native cache view to the measured scope}

Use \texttt{worked/cache-ordinary.asm} for the large ordinary observation and
\texttt{worked/cache-tiled.asm} for its blocked comparison. Each contains the
same deterministic signed inputs, shape $16,16,16,4$ and original allocation
sizes. Its header selects sixteen four-word lines with fully associative LRU
placement. Start a fresh Debug session and open Data Cache Simulator for that
execution. At \texttt{matmul\_call}, the arguments are ready, but the caller has
already read method: the launch view has one access and one miss. Keep that
state for Learn's reference launch, then resume to \texttt{matmul\_done}.
Repeat in a fresh launch of the other complete source at the same boundary.

The expected launch access counts are 8745 and 12841, with full misses 5388 and
1068. The main chapter's cold function table has one fewer access in each case.
The method read shares a line later used by the descriptor in these observations,
so the extra read does not add another final miss. That equality is a consequence
of the actual address trace and model, not a general rule for caller work.

\begin{figure}[htbp]
\centering
\begin{tikzpicture}[alt={Two actual light-theme Data Cache Simulator views at matmul_done. Both use sixteen four-word lines, fully associative LRU placement and 256-byte capacity. The ordinary launch has 8745 accesses, 3357 hits and 5388 misses. The blocked launch with tile edge four has 12841 accesses, 11773 hits and 1068 misses. Native boxes identify the counter rows. Both intervals include the caller method-word read.}]
\node[anchor=north west,inner sep=0,font=\small\bfseries] (atitle) at (0,0) {1. Ordinary product};
\node[anchor=north west,inner sep=0] (a) at ([yshift=-2mm]atitle.south west) {\MARSRevisionCrop{revision/captures/measurement/matrix-naive.png}{0.42\linewidth}{0.640625}{0.318888}{0.039062}{0.255556}};
\node[anchor=north west,inner sep=0,font=\small\bfseries] (btitle) at ([xshift=10mm]atitle.north east -| a.north east) {2. Blocked, $T=4$};
\node[anchor=north west,inner sep=0] (b) at (btitle.south west |- a.north west) {\MARSRevisionCrop{revision/captures/measurement/matrix-blocked.png}{0.42\linewidth}{0.640625}{0.318888}{0.039062}{0.255556}};
\foreach \panel in {a,b} {
 \begin{scope}[shift={(\panel.south west)},x={($(\panel.south east)-(\panel.south west)$)},y={($(\panel.north west)-(\panel.south west)$)}]
 \draw[black,line width=.7pt] (.01,.08) rectangle (.99,.23);
 \end{scope}
}
\end{tikzpicture}
\caption{Actual light-theme Debug observations at the same return boundary. Both use the measurement tool's deterministic input pattern, shape $16,16,16,4$ and fixed stack placement. Blocking reduces misses while increasing notices. Boxes are native LaTeX annotations. These launch observations include one more access than the cold function measurements.}
\label{fig:revision-measurement-cache}
\end{figure}


Before leaving each guided execution, continue to \texttt{results\_ready}.
The caller's store changes the ordinary totals to 8746 accesses and 5389 misses;
the blocked totals become 12842 and 1069. It is the same product measured over
a longer interval. Finish or stop Debug before the next launch. When ordinary
Run is requested, all five worked sources exit silently; their retained
\texttt{completed} counts are four, one, four, 256 and 256 for ordinary, wide,
blocked and the two large sources.

For a separate function-only experiment, start another fresh Debug session and
pause at \texttt{matmul\_call}. Select that execution in Data Cache Simulator
and click Reset while it is stopped. Verify zero accesses, hits and misses before
resuming to \texttt{matmul\_done}. The ordinary interval now has 8744 accesses
and 5388 misses; the blocked interval has 12840 and 1068. The call instruction
has no data notice, so this reset interval matches function entry. Record the
reset and both boundaries with the counters. A reset after partial multiplication
would measure a suffix instead; do not combine it with the reference launch.

Change only tile edge to eight and compare again under each stated placement.
Use the measured table to explain the direct-mapped regression. Increasing tile
size changes the working set and loop overhead together; it is not evidence that
larger tiles always improve a program. Keep capacity and replacement policy fixed
when investigating placement.

\begin{samepage}
  Run the three original correctness checks against your working copy:
  \begin{verbatim}
java -jar /path/to/cli.jar grade \
  packages/022-blocked-matmul/checks/kernel.kts \
  /path/to/your/matrix-measure.asm
java -jar /path/to/cli.jar grade \
  packages/022-blocked-matmul/checks/program.kts \
  /path/to/your/matrix-measure.asm
java -jar /path/to/cli.jar grade \
  packages/022-blocked-matmul/checks/reuse.kts \
  /path/to/your/matrix-measure.asm
\end{verbatim}
  \par\end{samepage}
These accept equivalent schedules under both public entries. The worked blocked
schedule is an explanatory and measurement example; exact instruction/miss counts
are not general correctness criteria.

\section{Implement a matrix-vector reduction independently}

Open \texttt{independent/matrix-vector.asm} in the imported lesson and replace
\texttt{matvec}'s failing body. Its original standalone copy is also supplied in
the lab files.
Arguments are matrix, vector, output and shape pointers in \texttt{\$a0} through
\texttt{\$a3}. The shape contains only two words $M,K$. The packed matrix has
$M$ rows and $K$ word columns; the vector has $K$ word elements. Each output is
one exact signed doubleword dot product. Return $M$.

Derive matrix row displacement $4iK$, vector displacement $4q$ and output
displacement $8i$. The vector restarts for each new row. The default inputs give
50 and 122. For a zero reduction width, store one zero doubleword per output row;
for zero rows, read neither input. Both input ranges and the selected output can
start inside their allocations, so use the passed pointers.

Preserve full saved registers, frame pointer and stack pointer, input/shape words,
guards and inactive output. Read only selected caller inputs, the eight-byte shape
and initialized output. Private stack storage is allowed. The output starts with
arbitrary bytes and must be initialized during this call before any partial-sum
read. Produce no console I/O and return through the supplied continuation.

\noindent\begin{minipage}{\linewidth}
  Run the original public script on your completed source:
  \begin{verbatim}
java -jar /path/to/cli.jar grade \
  revision/lessons/022-measuring-programs/independent/matrix-vector.kts \
  /path/to/your/matrix-vector.asm
\end{verbatim}
\end{minipage}
\par\medskip

Learn checks the current source in two activities. The first retains all eight
original function input families and the complete default caller, adding a direct
default-input function case and an explicit positive $2^{31}$ case. Each of the
ten function cases calls twice through different continuations; the eleventh
case checks the caller's two complete sums and published count before its exit
selector. Each case has a cumulative 100000-instruction budget. The second
activity requires the same source to continue to silent normal exit. The original
standalone script above retains its nine checks.

Final bytes and access histories answer different questions. Native checking
observes the assembled caller span, including unnamed alignment padding, while
permitting selected inputs, descriptor and initialized output. A temporary store
outside selected output is still forbidden after restoration. A read of old
output still fails after a later repair produces the right sum. Every output
byte must be initialized again in each call; earlier setup and a previous call
do not count. Private stack frames and different valid reduction orders remain
available. Exact instruction, traffic and miss counts are not acceptance targets.

Both native checks must pass for the current source. Editing a worked or
independent file clears its execution credit; Continue returns to the first
activity that needs fresh evidence while retaining prediction history. Keep
extra fault experiments in disposable copies so the guided sources remain
reproducible.

Keep a row invariant, a zero-$K$ trace and a signed-width counterexample with your
passing results. After establishing correctness, propose a measurement interval
for your own implementation and identify which loads/stores belong to its
arithmetic versus its interface overhead. Its cache counts are observations to
explain, not another acceptance target to imitate.
